\SourceLeaf{170}cordanzen bei den angesetzten Summationen zu erklären, die hin und wieder eine oder ein Paar Einheiten betragen können.

\Needspace{210pt}
\begin{center}
\textit{Bilanzrechnung der Witwenkasse für 1. October 1851.}\par\smallskip
\footnotesize\setlength{\tabcolsep}{1pt}
\begin{minipage}[t]{.48\linewidth}
\centering
Schuld.\par
\begin{tabular}{l r}
&Thaler\\
\shortstack[l]{Capitalwerth des festen\\Theils der Pensionen}&\\
\shortstack[l]{1. für die jetzt vorhandenen\\Witwen . .}&28786\\
\shortstack[l]{2. fur Witwen und Waisen\\der jetzigen Genossen\\. . . . . . . . . . . . . . . . . . .}&58463\\
\shortstack[l]{3. für Witwen und Waisen\\der künftig beitretenden\\Genossen . . . . . . . . . . . .}&40712\\
\shortstack[l]{4. Capitalwerth des beweglichen\\Theils der Pensionen nach\\jetziger Normirung . .}&25715\\
\shortstack[l]{5. Capitalwerth der sonstigen\\Ausgaben . . .}&12907\\\cline{1-2}
\multicolumn{1}{r}{Summa Thaler}&166581
\end{tabular}
\end{minipage}\hfill\vrule\hfill
\begin{minipage}[t]{.48\linewidth}
\centering
Gut.\par
\begin{tabular}{l r}
&Thaler\\
Capitalwerth&\\
\shortstack[l]{6. des Ertrags der Apotheke\\. . . . . . . . .}&28571\\
\shortstack[l]{7. der Beiträge der jetzigen\\Genossen . . .}&8961\\
\shortstack[l]{8. der Beiträge aller künftig\\eintretenden Genossen\\. . . . . . . . . . . . . . . . . . .}&13596\\
\shortstack[l]{9. Geldvermögen der\\Witwenkasse . . . . . .}&122290\\\cline{1-2}
\multicolumn{1}{r}{Summa Thaler}&173419
\end{tabular}\par\medskip
Das Resultat der Bilanzrechnung ist also ein \textit{Überschuss} von 6837 Thalern.
\end{minipage}
\end{center}
\noindent\hspace*{\parindent}Die einzelnen Positionen der vorstehenden Bilanzrechnung begleite ich mit folgenden nähern Erörterungen.

Zu (1). Der Capitalwerth der Witwenpensionen für sämmtliche 15 jetzt vorhandene Witwen nach dem festen Bestandtheile (zu 200 Thaler fur jede) ist die Summe der jetzigen Capitalwerthe dieser Pensionen für jede einzelne Witwe nach ihrem Lebensalter berechnet, mit Berücksichtigung eventueller Waisenpensionen, wo minorenne Kinder jetzt vorhanden sind. Ich setze diese Werthe einzeln hieher: die Witwen sind numerirt nach der vollständigen Reihenfolge seit Stiftung der Anstalt.
\begin{center}
\footnotesize\setlength{\tabcolsep}{3pt}
\begin{tabular}{r|l|r l||r|l|r l||r|l|r l}
31&\textsc{Sch.}&1066&Thaler&55&H.&1834&Thaler&65&H.&2003&Thaler\\
44&F.&1622&&58&G.&2730&&66&H.&1172&\\
50&H.&2849&&59&W.&1862&&68&M.&2923&\\
52&P.&1479&&62&B.&1003&&69&D.&2897&\\
53&M.&1812&&63&G.&1309&&70&L.&2224&
\end{tabular}
\end{center}
\noindent\hspace*{\parindent}Die Zahlen (2) und (7) sind durch folgendes Verfahren ermittelt, dessen Rechtfertigung in der Denkschrift von 1845 zu finden ist. Unter den 53 Mitgliedern, welche gegenwärtig die Genossenschaft ausmachen, sind zur Zeit 47 verheirathet. Für jedes derselben ist, nach Massgabe des Alters des Mannes und der Frau, der jetzige Capitalwerth sowohl der vom erstern zu leistenden Beiträge (zu 15 Thaler jährlich), als der der letztern im Fall des Überlebens zu Theil werdenden Witwenpension nach ihrem festen Theile (zu 200 Thaler) berechnet. So verstanden, ergibt sich die Summe der Beiträge zu 7947 Thaler, die Summe der Pensionen zu 45364 Thaler. Um die 6 jetzt unverheiratheten Mitglieder mit zu berücksichtigen, werden diese Zahlen mit \(\tfrac{53}{47}\) multiplicirt, woraus die Beiträge \(=8961\) Thaler, und die Pensionen \(=51155\) Thaler hervorgehen: erstere Zahl ist obige Position (7). Letztere, der Waisenpension wegen um \(\tfrac{1}{7}\) vergrössert, ergibt 58463 Thaler, die Position (2).

Die Rechtfertigung der Annahme des Bruches \(\tfrac{1}{7}\) für die Waisenpensionen, anstatt des 1845 angewandten Bruches \(\tfrac{1}{6}\), wird bei der Nachweisung der Positionen (3) und (8) gegeben werden.

Ich setze noch die Resultate obiger Rechnung für die einzelnen 47 verheiratheten Mitglieder hieher. Die Numerirung ist die Reihenfolge des Eintritts in die Anstalt; die Zahlen der dritten Columne sind die Beiträge, die der vierten die Pensionen.
